\chapter{Intestinal Obstruction Symptoms}

\section*{Definition}
Intestinal obstruction is a mechanical or functional blockage of intestinal luminal flow, classified as small bowel obstruction (SBO, accounting for 80\% of all obstructions) or large bowel obstruction (LBO). Diagnosis is confirmed by clinical presentation and imaging showing dilated bowel proximal to the obstruction with decompression distally.

\section*{What Happens in Your Body}
Mechanical obstruction prevents distal passage of intestinal contents, causing proximal bowel distension from accumulated fluid and gas. Intraluminal pressure rises, compromising venous drainage, reducing arterial perfusion, and ultimately leading to ischemia, necrosis, perforation, and bacterial translocation.

\section*{Causes}
\begin{itemize}
  \item \textbf{Small bowel obstruction:} Adhesions (prior abdominal surgery, 60--70\%), incarcerated hernia, neoplasm (primary or metastatic), Crohn disease, volvulus, intussusception, gallstone ileus, foreign body
  \item \textbf{Large bowel obstruction:} Colorectal carcinoma (most common), diverticulitis (inflammatory stricture), volvulus (sigmoid, cecal), fecal impaction, Ogilvie syndrome (acute colonic pseudo-obstruction), hernia
  \item \textbf{Paralytic ileus:} Post-operative (most common), electrolyte disorders (hypokalemia, hyponatremia, hypercalcemia), peritonitis, retroperitoneal hemorrhage, medications (opioids, anticholinergics)
  \item \textbf{Pseudo-obstruction:} Chronic intestinal pseudo-obstruction (CIPO), connective tissue disease (scleroderma), Parkinson disease, diabetes mellitus
\end{itemize}

\section*{When to See a Doctor}
See your doctor if:
\begin{itemize}
  \item Colicky abdominal pain with nausea, vomiting, absolute constipation (no flatus or stool for 24 hours)
  \item Abdominal distension, high-pitched bowel sounds (tinkling), or visible peristalsis
  \item Bilious or fecaloid vomiting
  \item Fever, tachycardia, or peritonitis (rigidity, guarding, rebound tenderness) suggesting strangulation
  \item Prior abdominal surgery with acute onset of obstruction symptoms
\end{itemize}

\section*{Related Symptoms}
\begin{itemize}
  \item Abdominal pain (colicky, intermittent), nausea, vomiting (bilious, then feculent), constipation, obstipation, abdominal distension, tympany, high-pitched bowel sounds
\end{itemize}
\textbf{Important:} Continuous severe pain with fever and peritonitis suggests strangulation (ischemic bowel) and requires urgent surgical exploration.

\section*{Self-Care / Home Management}
\begin{itemize}
  \item Nothing by mouth (NPO) immediately if obstruction is suspected
  \item Do not take laxatives or stool softeners (may increase intraluminal pressure and perforation risk)
  \item Seek emergency medical evaluation rather than attempting home management
  \item If partial obstruction known, a clear liquid diet may be tolerated under medical supervision
  \item Elevate the head of bed to reduce aspiration risk if vomiting occurs
\end{itemize}

\section*{How Your Doctor Will Evaluate You}
\begin{itemize}
  \item Abdominal X-ray (upright and supine): dilated bowel loops, air-fluid levels, absent distal air (step-ladder pattern in SBO, haustral pattern in LBO)
  \item CT abdomen/pelvis with IV and oral contrast: gold standard for identifying obstruction site, cause, and complications (ischemia, perforation)
  \item Abdominal ultrasound to rule out free fluid or assess for intussusception in children
  \item Laboratory: CBC (leukocytosis, left shift), BMP (electrolyte abnormalities, prerenal azotemia), lactate (elevated in ischemia)
  \item Water-soluble contrast study for evaluating partial SBO and potential for non-operative resolution
\end{itemize}

\section*{Treatment Options}
\begin{itemize}
  \item NPO, nasogastric tube decompression, IV fluid resuscitation, electrolyte correction
  \item Water-soluble contrast challenge: 50--150 mL diatrizoate meglumine (Gastrografin) may be both diagnostic and therapeutic
  \item Operative exploration for complete SBO, signs of strangulation, or failure of non-operative management (48--72 hours)
  \item Large bowel obstruction: decompression (colonoscopic stent or rectal tube) as bridge to elective resection, or emergency Hartmann procedure
  \item Sigmoid volvulus: endoscopic detorsion with rectal tube placement; elective sigmoidectomy if recurrent
\end{itemize}

\section*{References}

\begin{thebibliography}{99}
\bibitem{harrison-bowel-obstruction} Longo DL, et al. \emph{Harrison's Principles of Internal Medicine}. 21st ed. McGraw-Hill; 2022. Chapter 363: Intestinal Obstruction.
\bibitem{uptodate-sbo} Zielinski MD, et al. Management of small bowel obstruction. In: UpToDate, Post TW (Ed), UpToDate, Waltham, MA; 2025.
\bibitem{maung} Maung AA, et al. Small bowel obstruction: a clinical review. \emph{JAMA}. 2023;329(11):921-933.
\bibitem{catena} Catena F, et al. World Society of Emergency Surgery guidelines on intestinal obstruction. \emph{World J Emerg Surg}. 2024;19(1):15.
\bibitem{tenbroek} ten Broek RPG, et al. Adhesive small bowel obstruction: an updated review. \emph{Lancet Gastroenterol Hepatol}. 2023;8(4):356-367.
\bibitem{taylor} Taylor MR, et al. \emph{Sabiston Textbook of Surgery}. 21st ed. Elsevier; 2023. Chapter 48: Intestinal Obstruction.
\end{thebibliography}
